\chapter{Proposta do trabalho}

    \par Este capítulo tem como foco detalhar os processos para atingir os objetivo definidos no inicio do presente trabalho, a proposta consiste na refatoração de uma API REST, que originalmente faz o uso framework Laravel do PHP, para uma nova estrutura que segue os princípios da Arquitetura Limpa, objetivando demonstrar de maneira prática como realizar essa aplicação de arquitetura.
    
    \par O critério para a seleção dessa API REST, é que seja funcional e que resolva um problema específico e que essa API REST adote os preceitos do MVC do laravel de forma pura, assim, permitindo que o processo de refatoração para a Arquitetura Limpa e a comparação sejam focados e didáticos, ilustrando de forma muito clara a transição entre os dois modelos.

    \par Após essa escolha, é necessário um embasamento teórico sobre o framework Laravel e sobre sua estrutura MVC, seguindo sua documentação, a fim de justificar as alterações propostas na etapa de refatoração. Tendo essa base teórica, é possível implementar a Arquitetura Limpa reorganizando a estrutura da API REST, com o objetivo de transformar o MVC em um acessório da arquitetura, usado para orquestrar a comunicação do usuário com a a API, bem como fazer com que os detalhes do framework sejam mais desacoplados das regras de negócio.

    \par O esperado é um guia prático que documente o processo e demonstre de forma concreta como aplicar a Arquitetura Limpa em uma API REST feita Laravel com padrão MVC. Para que ao final seja conduzida uma análise comparativa entre a versão inicial e a versão refatorada da API REST selecionada, por meio de diagramas, exemplos de código em PHP e uma análise qualitativa geral do processo, assim cumprindo o objetivo final do trabalho.

    \par No Quadro \ref{quadro:cronograma}, tem-se o cronograma a ser seguido para entrega das atividades propostas que visam atingir os objetivos deste trabalho.

\captionsetup{labelsep=endash, justification=centering}

\begin{quadro}[H]
    \centering
    \caption{Cronograma TCC 2}
    \label{quadro:cronograma}
    \begin{tabularx}{\textwidth}{|p{4.4cm}|>{\centering\arraybackslash}X|>{\centering\arraybackslash}X|>{\centering\arraybackslash}X|>{\centering\arraybackslash}X|>{\centering\arraybackslash}X|}
        \hline
        \textbf{Atividades} &
        \textbf{Agosto} & 
        \textbf{Setembro} & 
        \textbf{Outubro} & 
        \textbf{Novembro} & 
        \textbf{Dezembro} \\
        \hline
        \hline
        Escolher uma API REST. & X & & & & \\
        \hline
        Análise na literatura sobre os conceitos do MVC no framework Laravel. & X &  & & & \\
        \hline
        Redesenho da API REST Laravel escolhida. & & X & X & X & \\
        \hline
        Comparação em termos de estrutura entre a arquitetura redesenhada com o modelo original da API REST. & & & X & X & \\
        \hline
        Entrega do TCC. & & & & X & \\
        \hline
        Apresentação do TCC. & & & & & X \\
        \hline
    \end{tabularx}
    
    \par\medskip % Adiciona um pequeno espaço vertical
    Fonte: Autor
\end{quadro}


    \par Este trabalho foi planejado e distribuído ao longo de cinco meses, a fase inicial concentra-se na fundamentação teórica e na seleção de uma API REST, a fase prática e de análise estende-se de setembro a novembro, começando com o redesenho da API REST escolhida, simultaneamente enquanto a etapa de redesenho da API REST é feita nos meses de outubro e novembro são documentadas as diferenças estruturais e as implicações do redesenho. O cronograma culmina com os marcos finais do projeto, a entrega do trabalho escrito em novembro e a apresentação formal dos resultados em dezembro.